\par\Needspace{6\baselineskip}\medskip
\hypertarget{native-217-article158}{}\pdfbookmark[2]{158.}{article-158}
\begin{center}158.\end{center}
The two preceding examples have afforded no occasion yet to use the method of Article~120: the successive hypotheses converged so rapidly that we could have stopped at the second, and the third differed scarcely perceptibly from the truth. We shall in fact always have this advantage and be able to dispense with a fourth hypothesis whenever the heliocentric motion is moderate and the three radius vectors are not too unequal, especially if the time intervals also differ little from one another. The further the conditions of the problem depart from these circumstances, however, the more the first assumed values of the quantities $P$, $Q$ will differ from the true ones and the more slowly the subsequent values will converge to the true ones. In such a case the first three hypotheses are still to be worked out as shown by the two preceding examples (with the sole difference that in the third hypothesis, just as in the first and second, we must compute the
\SourceLeaf{218}
quantities $\eta$, $\eta''$, $P'$, $Q'$, $X$, $Y$ instead of the elements themselves). The last values of $P'$, $Q'$ are then no longer to be taken as the new values of $P$, $Q$ in the fourth hypothesis. These new values are instead to be obtained by the method of Article~120 by combining the first three hypotheses. It will very rarely be necessary then to proceed to a fifth hypothesis according to the directions in Article~121. --- We shall now illustrate these computations too with an example, which will also show how widely applicable our method is.
